\section{Boundary Conditions}

Specification of domain boundary conditions for the particle concentration field $\phi$. A boundary condition must be specified for each of the six domain faces.

\begin{lstlisting}[language=C++]
BoundaryConditions {
    +x {
        <boundary condition parameters>
    }

    -x {
        <boundary condition parameters>
    }

    +y {
        <boundary condition parameters>
    }

    -y {
        <boundary condition parameters>
    }

    +z {
        <boundary condition parameters>
    }

    -z {
        <boundary condition parameters>
    }
}
\end{lstlisting}

\subsection*{Parameters}
The following parameter must be specified for the concentration field \lstinline!phi! within each boundary patch block.

phi
\vspace{\the\parameterListSpacing}
\begin{itemize}
    \item {\bf Description}: The boundary condition applied to the particle concentration field at this face.
    \item {\bf Required/Optional}: Required
    \item {\bf Type}: Named option.
    \item {\bf Options}:

        \begin{table}[H]
        \centering
        \begin{tabular}{|l|l|}
        \hline
        outflow    & \begin{tabular}[c]{@{}l@{}}Particles that reach this boundary are removed from the \\ simulation. Suitable for outflow and open boundaries. \end{tabular} \\ \hline
        reflection & \begin{tabular}[c]{@{}l@{}}Particles that reach this boundary are reflected (specular reflection) \\ back into the domain. Suitable for solid walls and symmetry planes. \end{tabular} \\ \hline
        periodic   & \begin{tabular}[c]{@{}l@{}}Particles that exit through this boundary re-enter through \\ the opposing face on the same axis. Must be applied to \\ both opposing faces simultaneously. \end{tabular} \\ \hline
        \end{tabular}
        \end{table}
    \item {\bf Example}: \lstinline!phi = outflow!
\end{itemize}
